% TOP_PROBLEM: {"problemId": "problem.p-versus-np", "canonicalTitle": "P versus NP", "releaseRank": 1, "primaryDomain": "theoretical_computer_science", "primaryDomainLabel": "Theoretical computer science", "displayStatus": "open", "releaseStatus": "open"}
% !TeX program = lualatex
% Definitions and sources retained from research_batch01.tex; research is in GPT_6_Astra_Ultra.
\documentclass[11pt]{article}
\usepackage[margin=25mm]{geometry}
\usepackage{fontspec}
\setmainfont{Latin Modern Roman}
\usepackage{amsmath,amssymb,mathtools}
\usepackage{unicode-math}
\setmathfont{Latin Modern Math}
\usepackage{enumitem,needspace}
\usepackage{xurl,hyperref}
\hypersetup{colorlinks=true,urlcolor=blue}
\setcounter{secnumdepth}{1}
\setlength{\parindent}{0pt}
\setlength{\parskip}{5pt}
\setlength{\emergencystretch}{3em}
\setlist[itemize]{leftmargin=3em,itemsep=5pt,topsep=4pt,parsep=0pt}
\allowdisplaybreaks
\newcommand{\N}{\mathbb{N}}
\newcommand{\Z}{\mathbb{Z}}
\newcommand{\Q}{\mathbb{Q}}
\newcommand{\R}{\mathbb{R}}
\newcommand{\C}{\mathbb{C}}
\newcommand{\F}{\mathbb{F}}
\newcommand{\A}{\mathbb{A}}
\newcommand{\PP}{\mathbb P}
\newcommand{\E}{\mathbb{E}}
\newcommand{\eps}{\varepsilon}
\newcommand{\dd}{\,\mathrm d}
\newcommand{\sref}[2]{\hyperlink{source:#1:#2}{[S#2]}}
\newcommand{\eref}[2]{\hyperlink{extra:#1:#2}{[E#2]}}
% Turn the next switch off for a shorter reading copy. The quotations preserve
% every exactTarget field, including qualifications omitted from a short summary.
\newif\ifshowcatalogue
\showcataloguetrue
\newcommand{\cataloguescope}[1]{\ifshowcatalogue
 \paragraph{Exact scope in the supplied catalogue.}
 \begin{quote}\small #1\end{quote}\fi}
\begin{document}
\setcounter{section}{0}
% ================================================================
% problemId: problem.p-versus-np
% source releaseRank: 1; source releaseStatus: open
% exactTarget SHA-256: 9ac81ed7e9ada6c6d9d77326bebb48362f3a3a5574e196911b79e5cc09752960
\clearpage
\section{\texorpdfstring{P versus NP}{P versus NP}}\label{1}
\subsection{Definitions and mathematical statement}
A language is a set $L\subseteq\{0,1\}^*$. Write $\mathsf P$ for languages decided by a deterministic Turing machine in time bounded by a polynomial in the input length. A language belongs to $\mathsf{NP}$ if there are a deterministic polynomial-time verifier $V$ and a polynomial $p$ such that
\[
 x\in L\quad\Longleftrightarrow\quad \exists y\in\{0,1\}^{\le p(|x|)}\;V(x,y)=1.
\]
The polynomial and machine may depend on $L$, but not on its individual inputs. The question is whether $\mathsf P=\mathsf{NP}$. A negative answer requires a language in $\mathsf{NP}$ outside $\mathsf P$, not merely a slow algorithm for one problem.
\par\smallskip\textit{Formulation and concept references:} \sref{1}{1}, \eref{1}{1}.
\cataloguescope{Determine whether P = NP, that is, whether every decision problem whose solutions can be verified in polynomial time by a deterministic Turing machine can also be solved in polynomial time by a deterministic Turing machine.}
\subsection{Short English statement}
Does an efficient way to check a proposed yes-answer always imply an efficient way to decide whether the answer is yes?
\subsection{Sources}
\begin{itemize}\raggedright
\item[\textnormal{[S1]}]\hypertarget{source:1:1}{} Clay Mathematics Institute, The Millennium Prize Problems.
\url{https://www.claymath.org/library/monographs/MPPc.pdf}.
{\small\textit{Catalogue formulation source.}}
\item[\textnormal{[S2]}]\hypertarget{source:1:2}{} P vs NP.
\url{https://www.claymath.org/millennium/p-vs-np/}.
{\small\textit{Catalogue research/status reference; not a proof endorsement.}}
\item[\textnormal{[S3]}]\hypertarget{source:1:3}{} Constructive solvability and the P versus NP problem — Arne Hole.
\url{https://arxiv.org/abs/2406.16843}.
{\small\textit{Catalogue research/status reference; not a proof endorsement.}}
\item[\textnormal{[S4]}]\hypertarget{source:1:4}{} PNP Labs — P versus NP formal reconstruction.
\url{https://pnplabs.com.au/}.
{\small\textit{Catalogue research/status reference; not a proof endorsement.}}
\item[\textnormal{[E1]}]\hypertarget{extra:1:1}{} Stephen Cook, The P versus NP Problem, Clay Mathematics Institute, Section 1 and Appendix.
\url{https://www.claymath.org/wp-content/uploads/2022/06/pvsnp.pdf}.
{\small\textit{Added primary source: Definitions and formulation. Consulted 24 September 2026.}}

% BEGIN ADDED RESEARCH SOURCES -- batch 1, problem 1
\item[\textnormal{[E2]}]\hypertarget{extra:1:2}{} Theodore Baker, John Gill and
Robert Solovay, \emph{Relativizations of the $\mathcal P=?\mathcal{NP}$ Question},
SIAM Journal on Computing 4(4) (1975), 431--442, oracle equality and separation
results. \url{https://doi.org/10.1137/0204037}.
{\small\textit{Consulted 24 September 2026: publisher abstract; used only for
the stated oracle obstruction, not a new proof audit.}}
\item[\textnormal{[E3]}]\hypertarget{extra:1:3}{} Alexander A. Razborov and Steven
Rudich, \emph{Natural Proofs}, Journal of Computer and System Sciences 55(1)
(1997), 24--35; author manuscript dated 30 November 1996, abstract and the
natural-property/hardness framework.
\url{https://www.cs.umd.edu/~gasarch/BLOGPAPERS/natural.pdf}.
{\small\textit{Consulted 24 September 2026. The pseudorandomness assumption is
retained; no unconditional barrier is asserted.}}
\item[\textnormal{[E4]}]\hypertarget{extra:1:4}{} Scott Aaronson and Avi Wigderson,
\emph{Algebrization: A New Barrier in Complexity Theory}, ACM Transactions on
Computation Theory 1(1) (2009), Article 2; author manuscript, Sections 1--2
and 5, especially Theorems 5.1 and 5.3.
\url{https://www.scottaaronson.com/papers/alg.pdf}.
{\small\textit{Consulted 24 September 2026: definitions and barrier statements.}}
\item[\textnormal{[E5]}]\hypertarget{extra:1:5}{} Ryan Williams,
\emph{Nonuniform ACC Circuit Lower Bounds}, Journal of the ACM 61(1) (2014),
Article 2, 32 pages; author manuscript, Theorem 1.1 and Section 1.1.
\url{https://people.csail.mit.edu/rrw/acc-lbs.pdf}.
\url{https://doi.org/10.1145/2559903}.
{\small\textit{Consulted 24 September 2026.}}
\item[\textnormal{[E6]}]\hypertarget{extra:1:6}{} Lijie Chen, Avishay Tal and
Yichuan Wang, \emph{Super-quadratic Lower Bounds for Depth-2 Linear Threshold
Circuits}, ECCC TR26-039, 15 March 2026, abstract and reported
$\mathsf E^{\mathsf{NP}}$ threshold-circuit bound.
\url{https://eccc.weizmann.ac.il/report/2026/039/}.
{\small\textit{Consulted 24 September 2026: report page and abstract. A
recent manuscript result, not an independent verification of its full proof.}}
\item[\textnormal{[E7]}]\hypertarget{extra:1:7}{} Adnan Darwiche and Pierre
Marquis, \emph{A Knowledge Compilation Map}, Journal of Artificial Intelligence
Research 17 (2002), 229--264, Sections 1--2 and the separation of succinctness,
queries and transformations.
\url{https://arxiv.org/abs/1106.1819}.
\url{https://doi.org/10.1613/jair.989}.
{\small\textit{Consulted 24 September 2026: article introduction and language
definitions. Used as methodological context, not as the source of Lemma 1.1.}}
% END ADDED RESEARCH SOURCES -- batch 1, problem 1
\end{itemize}

\end{document}
